\section{Failure Modes}

Lineage verification is most useful when failures are specific. The current
negative harnesses separate three adjacent problems:

\begin{center}
\small
\begin{tabularx}{\textwidth}{lXX}
\toprule
Negative case & Expected failure & Interpretation \\
\midrule
\texttt{bad-lineage} & Missing current reporting unit & A lineage event points to a unit that does not exist. \\
\texttt{stale-plan} & Missing plan unit summary & The RPLAN assignment expects a unit absent from the RCOUNT package. \\
\texttt{tampered-2028-source} & Source hash mismatch & The package source evidence changed after hashing. \\
\bottomrule
\end{tabularx}
\end{center}

These failures should not collapse into a generic "cross-cycle mismatch." A bad
lineage record asks the election package to fix its administrative history. A
stale plan asks the district aggregation layer to use the right unit universe.
A source hash mismatch asks the package publisher to explain why the evidence
bytes changed. The remedy differs because the failure layer differs.

The stale-plan negative is especially important for later district aggregation.
It fails before a verifier can safely publish district totals, because the plan
assignment refers to a reporting unit that the count package for that cycle does
not contain.
